\subsection{Um controle que preserva as janelas}
\label{sec:posicao}
Montamos os dois textos abaixo. O separador \(S\) é a frase
\textit{I have a few thoughts about this movie.}

\painel{\small\sffamily
  \textbf{A}\quad \tokenrel{S}\enspace
    \tokenrel[SoftTeal]{I loved the acting}\enspace\tokenrel{S}\enspace
    \tokenrel[SoftRust]{I hated the ending}\enspace\tokenrel{S}\par\smallskip
  \textbf{B}\quad \tokenrel{S}\enspace
    \tokenrel[SoftRust]{I hated the ending}\enspace\tokenrel{S}\enspace
    \tokenrel[SoftTeal]{I loved the acting}\enspace\tokenrel{S}}

O separador impede que uma janela de até cinco tokens alcance os dois
trechos. Verificamos que A e B contêm as mesmas janelas de IDs, com as
mesmas multiplicidades, para \(k=3,4,5\) nas duas tokenizações. Cada
filtro recebe o mesmo conjunto de ativações e devolve o mesmo máximo.
Os dois logits coincidiram em cada checkpoint. As duas versões expressam opiniões mistas.
\fonteatual{seminario/dados/pares_resultados.json}{Par e predições}.

\subsection{Dois picos por filtro separam o par. Eles ajudam no contraste?}
Congelamos embeddings e filtros das três CNNs por palavras. Para cada
largura, dividimos as janelas válidas em duas metades (o excedente fica
na primeira) e extraímos os máximos \(u\) e \(v\), com \(v=0\) se a segunda
metade estiver vazia. Reajustamos regressões logísticas em 6.144 resenhas
e escolhemos \(C\in\{10^{-3},10^{-2},10^{-1},1,10\}\) por CE na validação
de 1.536 resenhas. As duas versões de
600 valores usam os mesmos picos. Ordenar pelo valor, \((\min(u,v),\max(u,v))\),
descarta a metade de origem. Ordenar pela posição, \((u,v)\), a preserva.

\begin{tabularx}{\linewidth}{@{}Xrrr@{}}
\toprule
\tabelahead{Saída por filtro} & \tabelahead{Auditoria} &
  \tabelahead{Pares sintéticos} & \tabelahead{Pares editados}\\
\midrule
\LinhasPosicao
\bottomrule
\end{tabularx}
\legendinha{Tabela 8. Média de três sementes. A auditoria mede acurácia
por texto em 2.048 resenhas. Nas outras colunas, o par conta como acerto
se as duas versões forem classificadas corretamente.}

Os 256 pares sintéticos cruzam quatro moldes e dois conjuntos de oito
adjetivos. Dos 488 pares editados do IMDb Contrast Set \cite{gardner},
excluímos 144 com entrada de IDs já vista no treino ou na validação.
Oito dos 344 pares restantes ficam idênticos após o recorte, com rótulos
opostos. No par de controle, só a cabeça com posição deu margens
diferentes para A e B nas três sementes. Nos pares editados, ela acertou
39,15\%, ante 39,83\% com dois picos por valor. Nos dois moldes sintéticos
com veredito depois do contraste, todas as cabeças ficaram em 0 de 128
pares por semente.
\href{\projetourl/blob/24672b7c4fb090899c3517947f0132d1f6d6296d/seminario/dados/posicao_resultados.json}{Métricas e margens por semente}.
